\documentclass[11pt]{amsart}
\usepackage{amsmath,amssymb,amsthm,mathtools}
\newtheorem{theorem}{Theorem}[section]
\newtheorem{proposition}[theorem]{Proposition}
\newtheorem{definition}[theorem]{Definition}
\newtheorem{remark}[theorem]{Remark}
\title{Saturated Cartier D\'ecalage, Exceptional Recharge, and Source-Word Realization}
\author{Maximum-Likelihood Candidate Chapter}
\date{5 August 2026}
\begin{document}
\maketitle

\begin{abstract}
We propose a factorization-independent saturated transform associated with a
centre-enriched blowup word.  The bare modification morphism is insufficient,
but the final modification together with its total source-labelled exceptional
Cartier divisor and a finite power-torsion exponent determines a fully
saturated capsule.  The Berthelot--Ogus identity
$\eta_f\eta_g=\eta_{fg}$ collapses a finite divisor word to its product, while
an iterated cohomology formula removes all exceptional power torsion after a
finite Noetherian exponent.  The fully saturated capsule need not equal an
arbitrary iterated ordinary strict transform: later nonflat pullback can create
new torsion along an older exceptional component.  Their difference is a
finite exceptional-recharge packet.  We formulate the recharge-vanishing or
strict-descent alternative and its role in positive-characteristic
resolution.  The scheme-level theorem remains open.
\end{abstract}

\section{The minimal transform payload}
A strict transform depends on the chosen centre.  The identity morphism
obtained by blowing up the Cartier ideal $(z)$ on $\operatorname{Spec}k[z]$
annihilates the strict transform of $k[z]/(z)$, whereas the empty identity
modification does not.  The two presentations have different exceptional
Cartier divisors.  This suggests retaining
\[
 (\pi:Y\to X, E_\pi, N, Q_\pi,\text{source lineage})
\]
rather than the morphism alone, where $E_\pi$ is the total source-labelled
exceptional divisor, $N$ a torsion exponent, and $Q_\pi$ a saturated derived
transform.

\section{Cartier d\'ecalage}
Let $A$ be a ring and let $f\in A$ be a nonzerodivisor.  For a complex of
$f$-torsion-free modules, Berthelot--Ogus define $\eta_f$.  Its derived functor
satisfies
\[
 H^i(L\eta_fM)\cong f^i\bigl(H^i(M)/H^i(M)[f]\bigr).
\]
For nonzerodivisors $f,g$ one has
\[
 \eta_f\eta_gM=\eta_{fg}M.
\]
Induction gives the natural formula
\[
 H^i((L\eta_f)^NM)
 \cong f^{Ni}\bigl(H^i(M)/H^i(M)[f^N]\bigr).
\]
Indeed, if $Q_N=H^i(M)/H^i(M)[f^N]$, then the $f$-torsion of $Q_N$ is
$H^i(M)[f^{N+1}]/H^i(M)[f^N]$; one more application of the one-step formula
proves the assertion.  For degree zero this is the quotient by $f^N$-torsion.

\section{Finite saturation}
Let $M$ be coherent on a quasi-compact Noetherian scheme.  The ascending chain
\[
 M[f]\subset M[f^2]\subset\cdots
\]
stabilizes.  For a bounded perfect packet with finitely many coherent
cohomology sheaves, one exponent $N$ dominates all degrees and all members of a
finite portfolio.  We call
\[
 \operatorname{SCD}_{E,N}(M)=(L\eta_{\mathcal O(-E)})^NM
\]
the saturated Cartier-d\'ecalage capsule, with the predictable degree twists
retained in its data.

At the module level, if
\[
 \operatorname{Sat}_f(P)=\{x:f^nx\in P\text{ for some }n\},
\]
then
\[
 \operatorname{Sat}_{fg}(P)
 =\operatorname{Sat}_g(\operatorname{Sat}_f(P)).
\]
Hence a finite Cartier word is equivalent, at degree zero, to saturation by the
product equation and is invariant under permutation.

\section{One-step strict transform and the Tor firewall}
On an affine blowup chart with exceptional equation $a$, the ordinary strict
transform of a module $F$ is
\[
 (F\otimes B)/(a\text{-power torsion}).
\]
If $K=F\otimes_A^{\mathbf L}B$, then $H^0(K)=F\otimes_A B$ because derived
pullback is right $t$-exact.  Negative Tor remains in negative degrees and does
not enter degree zero.  Therefore, if $N$ dominates the $a$-power-torsion
exponent,
\[
 H^0((L\eta_a)^NK)
 \cong (F\otimes B)/(a\text{-power torsion}),
\]
which is the one-chart strict transform.  This is the degree-zero Tor
firewall.  It does not by itself compare a d\'ecalage already performed on one
model with a later nonflat pullback.

\section{The exceptional-recharge correction}
Let $F^{\mathrm{iter}}$ denote the actual iterated strict transform of $F$ along
a centre word $\pi$, and let $E_\pi$ be the total source exceptional divisor on
the final model.  Define
\[
 \operatorname{Recharge}_{E_\pi}(F^{\mathrm{iter}})
 =H^0_{E_\pi}(F^{\mathrm{iter}})
 =\bigcup_nF^{\mathrm{iter}}[I_{E_\pi}^n].
\]
The fully saturated transform is
\[
 F^{\mathrm{sat}}
 =F^{\mathrm{iter}}/
   \operatorname{Recharge}_{E_\pi}(F^{\mathrm{iter}}).
\]
A later nonflat pullback may create torsion along an older exceptional
component.  The subsequent ordinary strict-transform step need only remove
torsion along its new exceptional equation, so $F^{\mathrm{iter}}$ need not be
fully saturated along the total divisor.  Thus
\[
 F^{\mathrm{iter}}\cong F^{\mathrm{sat}}
 \quad\Longleftrightarrow\quad
 \operatorname{Recharge}_{E_\pi}(F^{\mathrm{iter}})=0.
\]
The total-divisor d\'ecalage capsule computes $F^{\mathrm{sat}}$, not
unconditionally $F^{\mathrm{iter}}$.

\section{Good triples and common refinements}
For a perfect complex $P$ and an effective Cartier divisor $E$, the canonical
blowing-up-complexes construction principalizes finitely many intrinsic ideals
and produces a universal good triple.  In that chamber $L\eta_E P$ remains
perfect and commutes with arbitrary pullback.  We propose to apply this
construction to the direct sum of the finite source comparison complexes.

After functorial flatification of a maximal carrier, the final owner transforms
are finitely presented and flat.  Their degree-zero exceptional torsion is
zero.  The good-triple packet controls the comparison complex under the
nonflat common-refinement morphism; the remaining project-specific assertion is
that the actual iterated word has zero recharge.

\section{Passive and active stopping rules}
For passive modules, apply enough d\'ecalage layers to remove all exceptional
power torsion.  For a marked contact equation, apply only the number of layers
prescribed by the controlled-transform mark.  The example $A/(u^m)$ gives the
local shadow
\[
 m\longmapsto m-1.
\]
The unconsumed layer count is part of the contact/debt rank.  Full passive
saturation and controlled active division are therefore two stopping rules for
one Cartier-layer calculus, not one transform semantics.

\begin{theorem}[Saturated Cartier d\'ecalage with recharge defect, target]
Let $\pi:Y\to X$ be a finite centre-enriched ordinary blowup word, let
$E_\pi$ be its total source-labelled exceptional Cartier divisor, and let $P$
be a perfect finite source packet with coherent cohomology.  There exists a
finite exponent $N$ such that the normalized iterated d\'ecalage
$(L\eta_{E_\pi})^NL\pi^*P$ has cohomology obtained by quotienting every
cohomology sheaf by its full exceptional power torsion.  It is independent of
the ordering and parenthesization of the divisor components, commutes with flat
base change, and, after canonical good-triple preparation, commutes with
arbitrary pullback.

For a degree-zero owner module, the actual iterated strict transform admits a
canonical comparison with the saturated capsule.  The kernel is the finite
exceptional-recharge module.  The comparison is an isomorphism exactly in the
no-recharge chamber.  Joint legality and the Regular-Flag Flat-Lift theorem
force recharge to vanish along legal centre words; a nonzero recharge packet
instead produces strict exceptional support/Fitting descent.  For a marked
contact packet, one prescribed layer is the controlled exceptional division.
All comparisons are compatible with charts, overlaps, owners, sources,
boundary, and debt.
\end{theorem}

\begin{remark}
This theorem is not proved.  The iterated cohomology formula has a standard
inductive proof; the module saturation product law, recharge coverage, and rank
arithmetic have been written in Lean.  The scheme actual-to-saturated
comparison, good-triple source packet, recharge vanishing/descent, and
hereditary interfaces remain open.
\end{remark}

\section{Publication architecture}
The current maximum-likelihood series consists of six papers totaling $564$
dense Annals/AMS-equivalent pages.  The present chapter is the transform core
of Paper III, \emph{Functorial Flatification, Saturated Cartier D\'ecalage,
Recharge Defects, and Hereditary Reentry}.

\end{document}
